% Figure from MATH 451 notes, section 23. Compile with the notes preamble.
\begin{tikzpicture}
\begin{axis}[nfig, width=0.52\textwidth, height=0.44\textwidth, clip=false,
  xmin=0, xmax=1.12, ymin=0, ymax=1.12,
  xtick={0,1}, ytick={1}, xlabel={}]
  \addplot[figB!30, thick, domain=0:1] {x} node[pos=0.55, above left, font=\scriptsize, figB!60] {$x$};
  \addplot[figB!45, thick, domain=0:1, samples=80] {x^2};
  \addplot[figB!60, thick, domain=0:1, samples=80] {x^4};
  \addplot[figB!75, thick, domain=0:1, samples=100] {x^8};
  \addplot[figB!90, thick, domain=0:1, samples=120] {x^16};
  \addplot[figB, thick, domain=0:1, samples=160] {x^32};
  \node[font=\scriptsize, figB] at (axis cs:0.73,0.40) {$x^{2}$};
  \node[font=\scriptsize, figB, anchor=west] at (axis cs:1.0,0.40) {$x^{32}$};
  \draw[figR, very thick] (axis cs:0,0) -- (axis cs:0.985,0);
  \addplot[only marks, mark=*, mark size=1.8pt, mark options={fill=white, draw=figR, thick}] coordinates {(1,0)};
  \addplot[only marks, mark=*, mark size=1.8pt, figR] coordinates {(1,1)};
  \node[font=\scriptsize, figR, anchor=west] at (axis cs:1.03,1) {$f(1) = 1$};
  \node[font=\scriptsize, figR, anchor=north] at (axis cs:0.5,-0.02) {$f = 0$ on $[0,1)$};
\end{axis}
\end{tikzpicture}
